\rhead{PR402: Major Project}
\phantomsection
\section*{Major Project (PR402)}
\label{course:PR402}

\noindent \small{\hyperref[main:scheme]{$\leftarrow$ Back to Scheme}} \vspace{0.5cm}

\noindent \textbf{Credits:} 8 (0-0-8)

\subsection*{Course Directive}
\noindent The Major Project is the capstone academic requirement of the B.Tech programme, carried out in Semester 8. It represents the culmination of the student's four years of technical education and demands the independent application of domain knowledge, software engineering discipline, and research methodology to a substantial, well-defined problem. Every student must complete a Major Project regardless of internship status or industry employment.

\vspace{0.3cm}

\noindent\textbf{Project Allocation} \\
Projects are allocated by the Department Project Coordination Committee at the beginning of Semester 7 to allow adequate preparation time. Each student is assigned a faculty supervisor whose research or consultancy interests align with the student's major specialization track. Students may propose a self-identified project topic, subject to approval by their supervisor and the committee. Industry-sponsored projects and continuation of internship work are permitted provided a faculty co-supervisor is assigned and the project scope satisfies the academic depth requirements defined below.

\vspace{0.3cm}

\noindent\textbf{Scope and Depth Requirements} \\
The Major Project must demonstrate the following minimum technical characteristics:

\vspace{0.2cm}

\begin{itemize}
    \item A clearly stated problem with justification of its relevance to the student's specialization track and the broader technical landscape.
    \item A survey of existing approaches or related work demonstrating familiarity with the state of the art.
    \item An original technical contribution: this may be a novel system design, an implementation benchmarked against baselines, a dataset with analysis, a simulation study with new findings, or a research result supported by experimental evidence.
    \item A working software, hardware, or simulation artefact that is demonstrable at the time of evaluation.
    \item Quantitative evaluation: results must be supported by appropriate metrics, baselines, and statistical or empirical validation.
\end{itemize}

\newpage
\rhead{Curriculum Schema}